% ============================================================
\setpartpreamble{\partintro
  {Intervening when saving, connecting and executing, and reporting events properly.}
  {
    \item Chapter~\ref{ch:interceptors}: Interceptors
    \item Chapter~\ref{ch:diagnostics-and-logging}: Diagnostics and logging
  }}
\part{Run time}\label{part:V}

% ------------------------------------------------------------
\chapter{Interceptors}\label{ch:interceptors}

Interceptors are hooks into EF's run time: before and after commands, connections,
transactions, SaveChanges and materialization. They can observe, change values, suppress
operations or replace results.

\begin{datatable}{H{1.00}H{1.50}H{0.50}}
  \textbf{Interface} & \textbf{Events} & \textbf{Lifetime} \\
  \tablehead
  \code{ISave\allowbreak{}Changes\allowbreak{}Interceptor} & SavingChanges, SavedChanges,
  SaveChangesFailed, ThrowingConcurrencyException & Scoped \\
  \tablerule
  \code{IDb\allowbreak{}Connection\allowbreak{}Interceptor} & Opening, closing, creating the
  connection & Scoped \\
  \tablerule
  \code{IDb\allowbreak{}Command\allowbreak{}Interceptor} & Creating and executing commands,
  data readers & Scoped \\
  \tablerule
  \code{IDb\allowbreak{}Transaction\allowbreak{}Interceptor} & Begin, commit, rollback,
  savepoints & Scoped \\
  \tablerule
  \code{IIdentity\allowbreak{}Resolution\allowbreak{}Interceptor} & Conflicts when attaching
  already tracked entities; never for queries (chapter~\ref{ch:materialization-from-rows-to-objects}) & Scoped \\
  \tablerule
  \code{IMaterialization\allowbreak{}Interceptor} & Creating and initializing entities
  (chapter~\ref{ch:materialization-from-rows-to-objects}) & \textbf{Singleton} \\
  \tablerule
  \code{IQuery\allowbreak{}Expression\allowbreak{}Interceptor} & LINQ expression before
  compilation & \textbf{Singleton} \\
\end{datatable}

The two singleton interceptors implement \code{ISingleton\allowbreak{}Interceptor}. They are
used once per internal provider, not per context, and must therefore be stateless and
thread-safe.

\section{Two ways to register}

\begin{datatable}{H{0.96}H{1.20}H{0.84}}
  \textbf{Way} & \textbf{Properties} & \textbf{When} \\
  \tablehead
  \code{options\allowbreak{}Builder.\allowbreak{}Add\allowbreak{}Interceptors(\allowbreak{}.\allowbreak{}.\allowbreak{}.\allowbreak{})}
  & Instances from the application; can use app services; differ per context & Interceptors
  with app dependencies (tenant, user, clock) \\
  \tablerule
  \code{services.\allowbreak{}Add\allowbreak{}Scoped<\allowbreak{}IInterceptor,\allowbreak{}
  X>(\allowbreak{})} in \code{Apply\allowbreak{}Services} & Resolved from the internal
  provider (“injected interceptors”); have access to EF services, not to app services &
  Purely technical interceptors without app dependencies \\
\end{datatable}

EF combines both ways into one list. The sample uses the first, because both interceptors
need the application's tenant (chapter~\ref{ch:public-api-and-integration-into}).

\section{The SaveChanges interceptor}

It sets the tenant of new entities and prevents data of another tenant from being written or
reassigned. This complements RLS (defense in depth) and produces clearer errors than the
database. Chapter~\ref{ch:bulk-operations-and-the-limits} shows which paths bypass it:

\begin{sourcecode}[Interceptors/\allowbreak{}Tenant\allowbreak{}Save\allowbreak{}Changes\allowbreak{}Interceptor.\allowbreak{}cs (\allowbreak{}1/\allowbreak{}2)]{csharp}
public class TenantSaveChangesInterceptor : SaveChangesInterceptor
{
    private readonly ITenantProvider _tenants;

    public TenantSaveChangesInterceptor(ITenantProvider tenants)
    {
        ArgumentNullException.ThrowIfNull(tenants);

        _tenants = tenants;
    }

    public override InterceptionResult<int> SavingChanges(
        DbContextEventData eventData, InterceptionResult<int> result)
    {
        if (eventData.Context is { } context) ApplyTenant(context);
        return result;
    }

    public override ValueTask<InterceptionResult<int>> SavingChangesAsync(
        DbContextEventData eventData, InterceptionResult<int> result,
        CancellationToken cancellationToken = default)
    {
        if (eventData.Context is { } context) ApplyTenant(context);
        return ValueTask.FromResult(result);
    }
\end{sourcecode}

\begin{sourcecode}[Interceptors/\allowbreak{}Tenant\allowbreak{}Save\allowbreak{}Changes\allowbreak{}Interceptor.\allowbreak{}cs (\allowbreak{}2/\allowbreak{}2,\allowbreak{} simplified)]{csharp}
    internal void ApplyTenant(DbContext context)
    {
        // Which entities are affected is stored in the runtime annotation
        // (chapter "Property-level Fluent API and runtime annotations")
        var changes = context.ChangeTracker.Entries()
            .Where(e => e.State is EntityState.Added or EntityState.Modified
                                   or EntityState.Deleted)
            .Select(e => (Entry: e, Name: e.Metadata.GetTenantPropertyName()))
            .Where(x => x.Name is not null)
            .Select(x => (x.Entry, Property: x.Entry.Property(x.Name!)))
            .ToList();
        if (changes.Count == 0) return;

        var current = _tenants.CurrentTenantId
            ?? throw new TenancyException("No active tenant ...");

        foreach (var (entry, property) in changes)
        {
            var currentValue = (TenantId)property.CurrentValue!;
            var originalValue = (TenantId)property.OriginalValue!;
            switch (entry.State)
            {
                case EntityState.Added when currentValue.IsEmpty:
                    property.CurrentValue = current;
                    break;
                case EntityState.Added when currentValue != current:
                    throw Mismatch(entry, currentValue, current);
                // Only an actual change of the value counts, not the flag
                case EntityState.Modified
                    when property.IsModified && currentValue != originalValue:
                    throw new TenancyException("The tenant property is immutable.");
                case EntityState.Modified or EntityState.Deleted:
                    if (entry.State == EntityState.Modified)
                        property.IsModified = false;    // do not write the column
                    // Empty = unknown (shadow property): the database decides
                    if (!originalValue.IsEmpty && originalValue != current)
                        throw Mismatch(entry, originalValue, current);
                    break;
            }
        }
    }
}
\end{sourcecode}

Anyone who only works with asynchronous code easily overrides only one variant:

\begin{sourcecode}[Interceptor (\allowbreak{}WRONG)]{csharp}
public class TenantSaveChangesInterceptor : SaveChangesInterceptor
{
    public override ValueTask<InterceptionResult<int>> SavingChangesAsync(
        DbContextEventData eventData, InterceptionResult<int> result,
        CancellationToken cancellationToken = default) { /* apply the tenant */ }

    // SavingChanges is missing
}
\end{sourcecode}

\begin{failure}[What happens if only the asynchronous variant is implemented?]
  \textbf{Observable.} With \code{Save\allowbreak{}Changes\allowbreak{}Async} everything
  works. A synchronous \code{Save\allowbreak{}Changes}, for example in a console tool or an
  older code path, saves without a tenant; RLS rejects the row with \code{WITH CHECK}, a
  plugin without database protection silently writes it wrongly.

  \textbf{Cause.} Depending on the call, EF calls exactly one variant. The base class
  \code{Save\allowbreak{}Changes\allowbreak{}Interceptor} implements the other one as an
  empty method, so there is no error. The same applies to connection and command
  interceptors.

  \textbf{Solution.} Always override both variants and call the same logic, like
  \code{Apply\allowbreak{}Tenant} in the sample.

  \textbf{Test:}
  \code{Failure\allowbreak{}Case\allowbreak{}Tests.\allowbreak{}Interceptor\allowbreak{}\_with\allowbreak{}\_only\allowbreak{}\_the\allowbreak{}\_async\allowbreak{}\_variant\allowbreak{}\_…}.
\end{failure}

\section{Tracking flag, change of value and database protection}\label{sec:tracking-flag-versus-value}

Three things look alike when saving and mean different things:

\begin{itemize}
  \item \textbf{The tracking flag.} \code{Db\allowbreak{}Context.\allowbreak{}Update} and
    \code{Attach} of an untracked entity set \code{Is\allowbreak{}Modified} for \emph{every}
    property, the tenant property included. Original and current value are the same: the value
    the entity brings along. Up to 10.0 the interceptor already rejected this flag as a change of
    tenant, and an ordinary update from an API failed.
  \item \textbf{The actual change of value.} A row is only reassigned when the value differs from
    the original value. The interceptor rejects that, even when the new value is the active
    tenant. For an unchanged value it clears the flag, so that the \code{UPDATE} does not write
    the tenant column.
  \item \textbf{The state of the database.} For a tracked entity the original value comes from
    the database. For a detached entity it is only a claim of the caller. A forged entity with
    the own tenant and the ID of a foreign row therefore passes every check in the interceptor.
    Only the policy decides: the \code{UPDATE} does not see the foreign row, changes zero rows,
    and EF reports a \code{Db\allowbreak{}Update\allowbreak{}Concurrency\allowbreak{}Exception}.
    An empty original value (a shadow property the entity cannot carry) counts as unknown, not
    as foreign; here, too, the database decides.
\end{itemize}

The interceptor therefore only promises what it can know without accessing the database: no
change of tenant in the change tracker, no values known to be foreign, no change without an
active tenant. Whether a row really belongs to the tenant is checked by RLS alone.
\code{Save\allowbreak{}Changes\allowbreak{}Interceptor\allowbreak{}Tests} covers the cases
without a database, \code{Detached\allowbreak{}Update\allowbreak{}Tests} against PostgreSQL with
a restricted role: \code{Update} with an unchanged tenant, synchronous and asynchronous, the
forged foreign row ID and the reassignment of a tracked row.

The logic lives in an \code{internal} method \code{Apply\allowbreak{}Tenant}. Tests can call
it directly through \code{Internals\allowbreak{}Visible\allowbreak{}To} without needing a
database (chapter~\ref{ch:test-strategy}).

\section{The connection interceptor}

It sets the session setting that the policies and
\code{Current\allowbreak{}Tenant(\allowbreak{})} evaluate every time a connection is opened.
The flow (intercept opening, sync and async, fetch the tenant) is the same for every
database and therefore lives in an abstract base in the core package
(chapter~\ref{ch:what-belongs-in-the-core}). The PostgreSQL package only supplies the
command. It gets the setting name from the \emph{internal} provider via the context from the
event data
(\code{context.\allowbreak{}Get\allowbreak{}Service<\allowbreak{}T>(\allowbreak{})}), and
the tenant from the \emph{app} container:

\begin{sourcecode}[Interceptors/\allowbreak{}Tenant\allowbreak{}Session\allowbreak{}Interceptor.\allowbreak{}cs (\allowbreak{}core),\allowbreak{} Runtime/\allowbreak{}Tenant\allowbreak{}Connection\allowbreak{}Interceptor.\allowbreak{}cs]{csharp}
// Core package
public abstract class TenantSessionInterceptor : DbConnectionInterceptor
{
    protected TenantSessionInterceptor(ITenantProvider tenants)
    {
        ArgumentNullException.ThrowIfNull(tenants);

        Tenants = tenants;
    }

    protected ITenantProvider Tenants { get; }

    public override async Task ConnectionOpenedAsync(
        DbConnection connection, ConnectionEndEventData eventData,
        CancellationToken cancellationToken = default)
    {
        await using var command =
            CreateSessionCommand(connection, eventData.Context, Tenants.CurrentTenantId);
        await command.ExecuteNonQueryAsync(cancellationToken);
    }

    // ConnectionOpened (synchronous) analogous

    protected abstract DbCommand CreateSessionCommand(
        DbConnection connection, DbContext? context, TenantId? tenant);
}

// PostgreSQL package
public class TenantConnectionInterceptor : TenantSessionInterceptor
{
    public TenantConnectionInterceptor(ITenantProvider tenants)
        : base(tenants)
    {
    }

    protected override DbCommand CreateSessionCommand(
        DbConnection connection, DbContext? context, TenantId? tenant)
    {
        var settingName = context?.GetService<ITenancySingletonOptions>()
                              .Settings.SettingName
                          ?? TenancySettings.Default.SettingName;

        var command = connection.CreateCommand();
        command.CommandText = "SELECT set_config(@name, @value, false)";
        AddParameter(command, "name", settingName);
        // Always set, even empty: no stale value remains
        AddParameter(command, "value", tenant?.ToString() ?? string.Empty);
        return command;
    }
}
\end{sourcecode}

\begin{note}[Empty string instead of NULL]
  Npgsql resets the session state when a connection returns to the pool. After that,
  \code{current\_\allowbreak{}setting(\allowbreak{}'app.\allowbreak{}tenant\_\allowbreak{}id',\allowbreak{}
  true)} returns an empty string instead of \code{NULL}. A \code{''::uuid} would abort every
  query. Policy and translator therefore use
  \code{nullif(\allowbreak{}.\allowbreak{}.\allowbreak{}.\allowbreak{},\allowbreak{} '')}:
  without a tenant the value is \code{NULL}, the comparison is never true and no rows are
  visible. That is the safe direction.
\end{note}

By default EF opens and closes the connection for every query. Every opening costs an extra
round trip here. With many small queries it pays off to open the connection once per request
explicitly (chapter~\ref{ch:plugin-performance}).

\section{Where the connection interceptor does not run}

The interceptor runs when EF opens a connection, and only then. Two situations bypass it,
and both are easy to create by accident:

\begin{datatable}{H{1.00}H{1.06}H{0.94}}
  \textbf{Situation} & \textbf{Consequence} & \textbf{Remedy} \\
  \tablehead
  An already open connection is passed in
  (\code{Use\allowbreak{}Npgsql(\allowbreak{}connection)},
  \code{Set\allowbreak{}Db\allowbreak{}Connection}) & EF never opens it, so the interceptor
  never runs. The session keeps whatever value it has: none, or that of a previous tenant &
  Pass connections closed and let EF open them; otherwise call
  \code{Database.\allowbreak{}Apply\allowbreak{}Tenant\allowbreak{}Session(\allowbreak{})}
  \\
  \tablerule
  The connection stays open (\code{Open\allowbreak{}Connection}, an open transaction) while
  the tenant changes & The new tenant only takes effect when the connection is opened again
  & One context, one tenant; if that is not possible, call
  \code{Database.\allowbreak{}Apply\allowbreak{}Tenant\allowbreak{}Session(\allowbreak{})}
  after the change \\
\end{datatable}

\begin{sourcecode}[tests/\allowbreak{}.\allowbreak{}.\allowbreak{}.\allowbreak{}Integration\allowbreak{}Tests/\allowbreak{}Connection\allowbreak{}Lifetime\allowbreak{}Tests.\allowbreak{}cs (\allowbreak{}shortened)]{csharp}
// Already open, with a session value left over from earlier work (tenant B)
await using var connection = new NpgsqlConnection(appConnectionString);
await connection.OpenAsync();
// ... SELECT set_config('app.tenant_id', <tenant B>, false) ...

// Meant for tenant A, but EF does not open the connection: the interceptor never runs
await using var db = CreateAppContext(connection, tenantA);
var orders = await db.Orders.ToListAsync();       // the rows of tenant B!
\end{sourcecode}

For these cases the plugin offers the explicit counterpart to the interceptor.
\code{Database.\allowbreak{}Apply\allowbreak{}Tenant\allowbreak{}Session(\allowbreak{})}
writes the current tenant into the session if the connection is open, and does nothing if it
is closed, because EF will then open it and the interceptor takes over. It reuses the
interceptor's own command, so it works for every provider package:

\begin{sourcecode}[Usage]{csharp}
await using var db = CreateAppContext(connection, tenantA);   // connection already open
await db.Database.ApplyTenantSessionAsync();                  // session now: tenant A
var orders = await db.Orders.ToListAsync();                   // the rows of tenant A
\end{sourcecode}

\begin{sourcecode}[Interceptors/\allowbreak{}Tenant\allowbreak{}Session\allowbreak{}Interceptor.\allowbreak{}cs (\allowbreak{}excerpt)]{csharp}
public virtual async Task<bool> ApplyToOpenConnectionAsync(
    DbContext context, CancellationToken cancellationToken = default)
{
    var connection = context.Database.GetDbConnection();
    if (connection.State != ConnectionState.Open)
        return false;                     // EF opens it later; the interceptor takes over

    await using var command =
        CreateSessionCommand(connection, context, Tenants.CurrentTenantId);
    // Inside an explicit transaction the command must be enlisted in it
    command.Transaction = context.Database.CurrentTransaction?.GetDbTransaction();
    await command.ExecuteNonQueryAsync(cancellationToken);
    return true;
}
\end{sourcecode}

The extension method \code{Apply\allowbreak{}Tenant\allowbreak{}Session} in the core package
finds the interceptor among the context's interceptors and fails with a clear message if
there is none. Two details matter: the method needs the context, not just the connection,
because the provider reads its settings from the internal service provider; that is why
\code{Create\allowbreak{}Session\allowbreak{}Command} receives the
\code{Db\allowbreak{}Context} instead of the event data. And the command must join a running
transaction: PostgreSQL is lenient here, SQL Server rejects a command that is not enlisted.
The SQL Server tests check exactly this case; without the line they fail.

\begin{important}[One context, one tenant]
  Isolation follows the connection, not the context. A context used for several tenants, or
  a connection that outlives a request, carries the tenant of the moment it was opened.
  \code{Connection\allowbreak{}Lifetime\allowbreak{}Tests} shows both cases against a real
  database, together with the two remedies: a closed connection that EF opens itself, and
  \code{Apply\allowbreak{}Tenant\allowbreak{}Session(\allowbreak{})}. A plugin whose hooks
  cannot cover every path should offer such an explicit API instead of only documenting the
  gap.
\end{important}

The gap can also be detected while loading: with
\code{Verify\allowbreak{}Materialized\allowbreak{}Tenants(\allowbreak{})} the plugin checks
every loaded entity against the context's tenant and throws instead of returning another
tenant's rows (chapter~\ref{ch:materialization-from-rows-to-objects}).

% ------------------------------------------------------------
\chapter{Diagnostics and logging}\label{ch:diagnostics-and-logging}

A professional plugin reports events through EF's logging system. Messages then end up in
the same channels users configure anyway (\code{ILogger}, \code{Log\allowbreak{}To},
\code{Diagnostic\allowbreak{}Source}), and warnings can be ignored, logged or turned into
errors via \code{Configure\allowbreak{}Warnings}.

\section{Event IDs}

\begin{sourcecode}[Diagnostics/\allowbreak{}Tenancy\allowbreak{}Event\allowbreak{}Id.\allowbreak{}cs (\allowbreak{}shortened)]{csharp}
public static class TenancyEventId
{
    public const int TenancyBaseId = 1_150_000;     // own, documented range

    private enum Id
    {
        TenantIsolationApplied = TenancyBaseId,
        TenantIsolationOptedOut,
        PoliciesWithoutRowLevelSecurity,
    }

    private static readonly string ModelPrefix = "Acme.EntityFrameworkCore.Tenancy.Model.";

    public static readonly EventId TenantIsolationApplied =
        Make(Id.TenantIsolationApplied);
    public static readonly EventId TenantIsolationOptedOut =
        Make(Id.TenantIsolationOptedOut);
    // ...

    private static EventId Make(Id id) => new((int)id, ModelPrefix + id);
}
\end{sourcecode}

EF and providers occupy low ranges. A plugin chooses its own high range and documents it.
Names follow EF's scheme “Assembly.Category.Event” so that prefix filters work.

\section{Event definitions}

An \code{Event\allowbreak{}Definition<\allowbreak{}T>} bundles event ID, log level, message
template and the warning behavior from the options. EF caches its definitions per internal
provider; the plugin does the same with its own class registered as a singleton
(chapter~\ref{ch:registering-services}):

\begin{sourcecode}[Diagnostics/\allowbreak{}Tenancy\allowbreak{}Logging\allowbreak{}Definitions.\allowbreak{}cs (\allowbreak{}excerpt)]{csharp}
public class TenancyLoggingDefinitions
{
    private EventDefinition<string>? _tenantIsolationOptedOut;

    public virtual EventDefinition<string> LogTenantIsolationOptedOut(
        IDiagnosticsLogger logger)
        => LazyInitializer.EnsureInitialized(ref _tenantIsolationOptedOut, () =>
            Create(logger, TenancyEventId.TenantIsolationOptedOut, LogLevel.Warning,
                nameof(TenancyEventId.TenantIsolationOptedOut),
                "Entity '{entityType}' implements ITenantOwned but is excluded from "
                + "tenant isolation."));

    private static EventDefinition<string> Create(
        IDiagnosticsLogger logger, EventId eventId, LogLevel level,
        string name, string message)
        => new(
            logger.Options,                       // contains the warning behavior
            eventId,
            level,
            "TenancyEventId." + name,             // appears in warning-as-error messages
            l => LoggerMessage.Define<string>(l, eventId, message));
}
\end{sourcecode}

\section{The logger extension}

The actual log method follows the pattern of EF's
\code{Core\allowbreak{}Logger\allowbreak{}Extensions} exactly: classic logging first, then
the event data for DiagnosticSource and \code{Log\allowbreak{}To}:

\begin{sourcecode}[Diagnostics/\allowbreak{}Tenancy\allowbreak{}Logger\allowbreak{}Extensions.\allowbreak{}cs (\allowbreak{}excerpt)]{csharp}
private static void LogEntityType(
    IDiagnosticsLogger<DbLoggerCategory.Model> diagnostics,
    EventDefinition<string> definition,
    IReadOnlyEntityType entityType)
{
    // 1. ILogger; throws with ConfigureWarnings(w => w.Throw(...))
    if (diagnostics.ShouldLog(definition))
        definition.Log(diagnostics, entityType.DisplayName());

    // 2. DiagnosticSource and LogTo (simple logging)
    if (diagnostics.NeedsEventData(definition, out var diagnosticSource, out var simpleLog))
    {
        var eventData = new EntityTypeEventData(definition, GenerateMessage, entityType);
        diagnostics.DispatchEventData(definition, eventData, diagnosticSource, simpleLog);
    }
}
\end{sourcecode}

Users can then turn the warning into an error:

\begin{sourcecode}[Usage]{csharp}
options.UseNpgsql(cs)
       .UseAcmeTenancy()
       .ConfigureWarnings(w => w.Throw(TenancyEventId.TenantIsolationOptedOut));

// -> InvalidOperationException while building the model:
//    "An error was generated for warning '...TenantIsolationOptedOut': Entity
//     'AuditEntry' ... This exception can be suppressed or logged by passing event ID
//     'TenancyEventId.TenantIsolationOptedOut' to the 'ConfigureWarnings' method ..."
\end{sourcecode}

\begin{tip}
  Security-relevant deviations belong in the log as a \textbf{warning}, not as a debug
  message. Teams with strict requirements escalate them to errors via
  \code{Configure\allowbreak{}Warnings}; others at least see them at startup. The plugin
  does not have to make this decision for everyone.
\end{tip}
